% Part: second-order-logic
% Chapter: sol-and-sets
% Section: comparing-sets

\documentclass[../../../include/open-logic-section]{subfiles}

\begin{document}

\olfileid{sol}{set}{cmp}

\olsection{د سټونو پرتله کول}

\begin{prop}
!!{formula} $\lforall[x][(X(x) \lif Y(x))]$ د فرعي سټ اړيکه
تعريفوي؛ يعنې $\Sat{M}{\lforall[x][(X(x) \lif Y(x))]}[s]$ هله او
يوازې هله چې $s(X) \subseteq s(Y)$۔
\end{prop}

\begin{prop}
!!{formula} $\lforall[x][(X(x) \liff Y(x))]$ د سټونو د يووالي
اړيکه تعريفوي؛ يعنې $\Sat{M}{\lforall[x][(X(x) \liff Y(x))]}[s]$
هله او يوازې هله چې $s(X) = s(Y)$۔
\end{prop}

\begin{prop}
!!{formula} $\lexists[x][X(x)]$ د ناتشوالي خاصيت تعريفوي؛ يعنې
$\Sat{M}{\lexists[x][X(x)]}[s]$ هله او يوازې هله چې $s(X) \neq
\emptyset$۔
\end{prop}

سټ~$X$ له سټ~$Y$ څخه لوے نۀ دے، $\cardle{X}{Y}$، هله او يوازې
هله چې !!a{injective} تابع $f\colon X \to Y$ وي۔ دا چې د تابع
يو پر يو والے او په~$X$ کښې د ورودي غړو د ارزښتونو په~$Y$
کښې شتون څرګندولے شو، د تعريف ساحې د فرعي سټونو د اندازې د
«لوے نۀ دے» اړيکه هم تعريفولے شو۔

\begin{prop}
لاندې فورمول
\[
\lexists[u][(\lforall[x][(X(x) \lif Y(u(x)))] \land \lforall[x][\lforall[y][(\eq[u(x)][u(y)] \lif
      \eq[x][y])]])]
\]
د «لوے نۀ دے» اړيکه تعريفوي۔
\end{prop}

دوه سټونه يو شان اندازه لري، يا «همشمېره» دي، $\cardeq{X}{Y}$،
هله او يوازې هله چې !!a{bijective} تابع~$f\colon X \to Y$ وي۔

\begin{prop}
لاندې فورمول
\begin{multline*}
  \lexists[u][(\lforall[x][(X(x) \lif Y(u(x)))] \land {}\\
    \lforall[x][\lforall[y][(\eq[u(x)][u(y)] \lif \eq[x][y])]]
  \land {} \\\lforall[y][(Y(y) \lif \lexists[x][(X(x)
      \land \eq[y][u(x)])])])]
\end{multline*}
د همشمېره‌والي اړيکه تعريفوي۔
\end{prop}

دا !!{formula}s به په ترتيب سره په $X \subseteq Y$، $X = Y$،
$X \neq \emptyset$، $\cardle{X}{Y}$ او $\cardeq{X}{Y}$ لنډوو۔
(دا لږ څه ګډوډوونکې کېدے شي، ځکه د سټونو~$X$ او~$Y$ په
غيررسمي خبرو کښې هم همدا نښې کاروو؛ خو دلته دا نښې د دويمې درجې
منطق د هغو !!{formula}s لنډې بڼې دي چې يوځايي اړيکيز متغيرونه
$X$ او~$Y$ پکښې راځي۔)

\begin{prop}
!!{sentence} $\lforall[X][\lforall[Y][((\cardle{X}{Y} \land
    \cardle{Y}{X}) \lif \cardeq{X}{Y})]]$ معتبره ده۔
\end{prop}

\begin{proof}
دا !!{sentence} په !!a{structure}~$\Struct{M}$ کښې هله صدق کوي
چې د هر دوو فرعي سټونو $X \subseteq \Domain{M}$ او
$Y \subseteq \Domain{M}$ دپاره، که $\cardle{X}{Y}$ او
$\cardle{Y}{X}$ وي، نو $\cardeq{X}{Y}$ هم وي۔ خو دا د
\emph{هر ډول} سټونو~$X$ او~$Y$ دپاره رښتيا ده؛ همدا د
شرودرـبرنشتاين قضيه ده۔
\end{proof}


\end{document}
